%% kapa: exam=1 round=35 year=2024 period=1 form=A subject=real_estate subject_ko=부동산학원론 seq=20 number=100
%% source: Q-Net 감정평가사 기출문제 (cst00309) artlSeq=5213229 file=raw/1008/5213229/1교시 문제지 원본.pdf
%% conversion: kordoc 4.24.51 -> tools/gen1.py (auto); review=proofread-sonnet; flags=none
\begin{question}{100}
대상물건에 관한 감정평가방법으로 옳지 않은 것은? (단, 주어진 조건에 한함)
\begin{choices}
  \item 주택으로 쓰는 층수가 4개 층으로 1개 동의 바닥면적의 합계가 700제곱미터인 건물에서 구분소유 부동산의 감정평가액은 합리적인 배분기준에 따라 토지가액과 건물가액으로 구분하여 표시할 수 있다.
  \item 주택으로 쓰는 층수가 3개 층으로 15세대가 거주할 수 있고 주택으로 쓰이는 바닥면적의 합계가 600제곱미터인 1개 동이며 구분소유가 아닌 건물의 감정평가는 토지와 건물을 일괄평가하는 것을 원칙으로 한다.
  \item 주택으로 쓰는 층수가 6개 층인 건물에서 구분소유 부동산의 감정평가는 거래사례비 교법으로 하는 것을 원칙으로 한다.
  \item 주택으로 쓰는 층수가 4개 층으로 1개 동의 바닥면적의 합계가 500제곱미터인 건물에서 구분소유 부동산의 감정평가는 토지와 건물을 일괄평가하는 것을 원칙으로 한다.
  \item 구분소유 부동산을 감정평가할 때에는 층별ㆍ위치별 효용요인을 반영하여야 한다.
\end{choices}
\end{question}
\answer{2}
